%! AP Statistics — Unit 2, Lesson 2.4 — Group Activity (Answer Key)
\documentclass[10pt]{article}
\usepackage{apstats-article}
\usepackage{apstats-key}

%=============================================================================
\begin{document}

\pageheader{Unit 2, Lesson 2.4}{Group Activity --- Answer Key}
\namepartnerperiod

\vspace{0.18in}

\begin{scenariobox}[Context]{sky}
\small
Your group has been given a real-world dataset. Your task is to construct a scatterplot,
describe the relationship using the DOFS framework, and draw a conclusion in context.
Each tier uses a different dataset --- compare results with other groups at the end.
\end{scenariobox}

\vspace{0.12in}

% ── TIER R ────────────────────────────────────────────────────────────────────
\begin{tcolorbox}[colback=white, colframe=black!40, boxrule=0.8pt, arc=2mm,
    left=4mm, right=4mm, top=3mm, bottom=3mm,
    title={\textbf{Tier R --- Remediate}}, fonttitle=\bfseries, halign title=center]
\small

\textbf{Dataset: Study time vs.\ quiz score (12 students)}

\vspace{6pt}
\renewcommand{\arraystretch}{1.4}
\begin{tabularx}{\linewidth}{@{} l *{12}{C{0.72cm}} @{}}
\rowcolor{sky}
\textbf{Study (hr)} & 0 & 0.5 & 1 & 1 & 1.5 & 2 & 2 & 2.5 & 3 & 3 & 3.5 & 4 \\
\textbf{Score}      & 55 & 60 & 62 & 65 & 68 & 72 & 75 & 78 & 80 & 84 & 86 & 90 \\
\end{tabularx}

\vspace{10pt}

\begin{enumerate}[leftmargin=1.6em, itemsep=0.18in]

  \item Plot each point on the grid below. Label both axes (variable name and units).

        \begin{center}
        \begin{tikzpicture}
          \draw[very thin, gray!35] (0,0) grid (8,6);
          \draw[->, thick, navy] (0,0) -- (8.4,0);
          \draw[->, thick, navy] (0,0) -- (0,6.5);
          \node[below=8pt] at (4,0) {\small \ans{Study time (hours)}};
          \node[left=8pt, rotate=90] at (0,3) {\small \ans{Quiz score (points)}};
        \end{tikzpicture}
        \end{center}

        \begin{teachernote}
        $x$-axis: Study time (0--4 hr), $y$-axis: Quiz score (55--90).
        Points should rise from left to right with little scatter --- strong positive linear.
        \end{teachernote}

  \item Write a complete DOFS description in context.\\[8pt]
        \ansline{There is a strong positive linear association between study time and quiz score}\\[2pt]
        \ansline{for these 12 students. As study time increases, quiz scores tend to increase.}\\[2pt]
        \ansline{There are no obvious outliers.}

  \item Does any point appear to be an outlier? If so, describe it in context and speculate
        on a possible reason.\\[8pt]
        \ansline{No obvious outliers in this dataset. Accept any well-reasoned response.}\\[2pt]
        \writeline

\end{enumerate}
\end{tcolorbox}

\vspace{0.12in}

% ── TIER A ────────────────────────────────────────────────────────────────────
\begin{tcolorbox}[colback=white, colframe=black!40, boxrule=0.8pt, arc=2mm,
    left=4mm, right=4mm, top=3mm, bottom=3mm,
    title={\textbf{Tier A --- Approaching Proficiency}}, fonttitle=\bfseries, halign title=center]
\small

\textbf{Dataset: Age of car (years) vs.\ resale price (\$1000s) --- 12 cars}

\vspace{6pt}
\renewcommand{\arraystretch}{1.4}
\begin{tabularx}{\linewidth}{@{} l *{12}{C{0.72cm}} @{}}
\rowcolor{sky}
\textbf{Age (yr)} & 1 & 2 & 3 & 3 & 4 & 5 & 5 & 6 & 7 & 8 & 10 & 12 \\
\textbf{Price}    & 28 & 25 & 22 & 20 & 18 & 16 & 14 & 13 & 11 & 8 & 6 & 3 \\
\end{tabularx}

\vspace{10pt}

\begin{enumerate}[leftmargin=1.6em, itemsep=0.18in]

  \item Identify the explanatory and response variables. Plot each point on the grid below.
        Label both axes completely.

        Explanatory: \ans{age of car (years)} \quad Response: \ans{resale price (\$1000s)}

        \begin{center}
        \begin{tikzpicture}
          \draw[very thin, gray!35] (0,0) grid (8,6);
          \draw[->, thick, navy] (0,0) -- (8.4,0);
          \draw[->, thick, navy] (0,0) -- (0,6.5);
          \node[below=8pt] at (4,0) {\small \ans{Age of car (years)}};
          \node[left=8pt, rotate=90] at (0,3) {\small \ans{Resale price (\$1000s)}};
        \end{tikzpicture}
        \end{center}

  \item Write a complete DOFS description in context.\\[8pt]
        \ansline{There is a strong negative linear association between age and resale price}\\[2pt]
        \ansline{for these 12 cars. As the age of the car increases, the resale price tends to decrease.}\\[2pt]
        \ansline{There are no obvious outliers among the original 12 data points.}

  \item A 50-year-old classic car has a resale price of \$45{,}000. Plot this point. Does it
        follow the overall pattern? Explain what kind of point this is and why it is unusual.\\[8pt]
        \ansline{No --- it is an outlier. The 50-year-old car has an extremely high price compared}\\[2pt]
        \ansline{to what the linear pattern would predict. Possible reason: it is a collector/vintage}\\[2pt]
        \ansline{car whose value increases with age rather than decreasing.}

\end{enumerate}
\end{tcolorbox}

\vspace{0.12in}

% ── TIER E ────────────────────────────────────────────────────────────────────
\begin{tcolorbox}[colback=white, colframe=black!40, boxrule=0.8pt, arc=2mm,
    left=4mm, right=4mm, top=3mm, bottom=3mm,
    title={\textbf{Tier E --- Extension}}, fonttitle=\bfseries, halign title=center]
\small

\textbf{Dataset: Distance from city center (miles) vs.\ average monthly rent (\$) --- 10 neighborhoods}

\vspace{6pt}
\renewcommand{\arraystretch}{1.4}
\begin{tabularx}{\linewidth}{@{} l *{10}{C{0.82cm}} @{}}
\rowcolor{sky}
\textbf{Distance} & 0.5 & 1 & 2 & 3 & 4 & 5 & 7 & 9 & 10 & 12 \\
\textbf{Rent}     & 3200 & 2900 & 2500 & 2200 & 2000 & 1800 & 1600 & 1400 & 1300 & 800 \\
\end{tabularx}

\vspace{10pt}

\begin{enumerate}[leftmargin=1.6em, itemsep=0.18in]

  \item Construct the scatterplot on the grid below, choosing appropriate axis scales.

        \begin{center}
        \begin{tikzpicture}
          \draw[very thin, gray!35] (0,0) grid (8,6);
          \draw[->, thick, navy] (0,0) -- (8.4,0);
          \draw[->, thick, navy] (0,0) -- (0,6.5);
          \node[below=8pt] at (4,0) {\small \ans{Distance from city center (miles)}};
          \node[left=8pt, rotate=90] at (0,3) {\small \ans{Average monthly rent (\$)}};
        \end{tikzpicture}
        \end{center}

  \item Write a complete DOFS description in context.\\[8pt]
        \ansline{There is a strong negative linear (or slightly curved) association between}\\[2pt]
        \ansline{distance from the city center and average monthly rent for these 10 neighborhoods.}\\[2pt]
        \ansline{As distance increases, rent tends to decrease. No obvious outliers.}

  \item The pattern looks roughly linear, but could also be slightly curved. Explain how
        you would decide which description is more accurate, and why the form matters for
        what we do next in the course.\\[8pt]
        \ansline{Look at whether the points follow a straight line closely or bend. If there is}\\[2pt]
        \ansline{a consistent curve, report ``curved.'' Form matters because correlation and least-squares}\\[2pt]
        \ansline{regression (Lessons 2.5--2.8) only apply to linear associations.}

  \item A neighborhood 3 miles from the center has a monthly rent of \$3{,}500. Add this
        point to your scatterplot. Describe it in context and speculate on why rent might
        be so high despite being far from the center.\\[8pt]
        \ansline{This point is an outlier --- much higher rent than predicted at 3 miles.}\\[2pt]
        \ansline{Possible reason: proximity to a desirable school, park, transit hub, or new development.}

\end{enumerate}
\end{tcolorbox}

\end{document}
